% s07-character.tex —— §5 The character sheaf property，原文 p.29–p.31
\srcpage{29}

\section{特征层性质}
\label{sec:character}

回忆 $\cF_\varphi$ 是一个 Weil 层，因此它带有从 $\Bun_{G,\ol{\FF}_q}$ 到 $\Bun_G$ 的 Frobenius 下降数据\yizhu{下降数据指把 $\Bun_{G,\ol{\FF}_q}$ 上的层连同 Frobenius 作用下的粘合同构一起下降到 $\Bun_G$ 上；Weil 层正是带有这种结构的层。}{原文 \S5, p.29}。
取 $\delta\in G(E)$，并把它看作 $G(L)$ 的元素\yizhu{$L=\widehat{E}^{\mathrm{un}}$ 是 $E$ 的最大非分歧扩域之完备化，$\sigma$ 为其 Frobenius，见 \S2.2。}{原文 \S2.2, p.19}。
态射 $\nu_\delta:\mathbb{D}\to G$ 定义在 $E$ 上，因为 Dieudonné--Manin 斜率分解定义在 $E$ 上（给定有限维 $E$-向量空间的自同构，\srcpageinl{30}按其特征多项式的根的绝对值，把该多项式分解为乘积）。
此外，$\delta$ 落在 $\nu_\delta$ 的中心化子中。
于是，若 $\delta$ 在 $G(E)$ 中椭圆，则它在 $G(L)$ 中是基本的。
我们由此得到群胚之间的函子
\[
[\{G(E)\}_{\mathrm{ell}}/\text{共轭}]
\longrightarrow
[G(L)_{\mathrm{basic}}/\sigma\text{-共轭}],
\]
以及一个映射
\[
\{G(E)\}_{\mathrm{ell}}\longrightarrow \mathrm{B}(G)_{\mathrm{basic}}.
\]

现设 $\delta$ 是椭圆的。
此时态射
\[
x_\delta:\Spa(\ol{\FF}_q)\longrightarrow \Bun_{G,\ol{\FF}_q}
\]
定义在 $\FF_q$ 上。
利用 $\cF_\varphi$ 的 Weil 层结构，$x_\delta^*\cF_\varphi$ 带有 Frobenius $\mathrm{Frob}$ 的作用。

通过同构
\[
[\Spa(\ol{\FF}_q)/J_\delta(E)]
\xrightarrow{\ \sim\ }
\Bun^{ss,\kappa(\delta)}_{G,\ol{\FF}_q},
\]
右边的 Frobenius 下降数据由左边的态射对
\[
(\mathrm{Frob},\sigma):(\Spa(\ol{\FF}_q),J_\delta(E))
\longrightarrow(\Spa(\ol{\FF}_q),J_\delta(E))
\]
给出，其中 $\sigma$ 被看作 $G(L)$ 的自同构，并限制为 $J_\delta(E)$ 的自同构。
但此时对 $g\in J_\delta(E)$，有 $g\delta g^{-\sigma}=\delta$（其中 $g\in G(L)$），于是
\[
g^\sigma=\delta^{-1}g\delta.
\]
由此可知，该 Frobenius 下降数据由
\[
(\mathrm{Frob},\mathrm{Int}_{\delta^{-1}}):(\Spa(\ol{\FF}_q),J_\delta(E))
\longrightarrow(\Spa(\ol{\FF}_q),J_\delta(E))
\]
给出。

于是，$\mathrm{Frob}$ 在 $\ol{\QQ}_\ell$-向量空间 $V=x_\delta^*\cF_\varphi$ 上的作用由某个自同构 $u\in\mathrm{GL}(V)$ 给出，它满足
\[
\forall g\in J_\delta(E),\qquad
u\circ\pi(g)=\rho(\delta)^{-1}\circ\pi(g)\circ\rho(\delta)\circ u,
\]
其中 $\pi:J_\delta(E)\to\mathrm{GL}(V)$。
特征层性质于是要求 $u=\pi(\delta)$\yizhuw{此处 $\rho(\delta)$ 在原文中并未定义：它应是把 $\pi$ 延拓到 $G(L)$ 上的表示 $\rho$ 在 $\delta$ 处的取值。译文照原式保留。}。

这里是对该条件的另一种表述。
存在分解
\[
\pi=\bigoplus_{\substack{\rho\in\widehat{S}_\varphi\\ \rho|_{Z(\widehat{G})^\Gamma}=\kappa(\delta)}}
\rho\otimes\pi_\rho,
\]
它与 $\mathrm{Frob}$ 的作用交换（按假设，Weil 下降数据与 $S_\varphi$ 的作用交换）。
既然按假设对每个 $\rho$，表示 $\pi_\rho$ 都不可约，就存在一族标量 $(\lambda_\rho)_\rho\subset\ol{\QQ}_\ell^\times$，使得 $\mathrm{Frob}$ 的作用由 $\bigoplus_\rho\lambda_\rho\pi(\delta)$ 给出\yizhuw{原文这里写的是 $\pi(\delta)$，而非 $\pi_\rho(\delta)$；照原样排印。}。
特征层性质于是要求：对一切 $\rho$，$\lambda_\rho=1$。

下面是由引入特征层性质得到的一个推论。

\textbf{特征层性质的推论：}\ 设 $T_\varphi$ 是 $G(E)$ 上联系于 Langlands 参数 $\varphi$ 的稳定分布。
$T_\varphi$ 在 $G(E)$ 的椭圆正则子集上的限制，由 Frobenius 迹函数
\[
\delta\longmapsto e(J_\delta)\,\mathrm{Tr}\bigl(\mathrm{Frob};x_\delta^*\cF_\varphi\bigr)
\]
经映射 $\{G(E)\}_{\mathrm{ell}}\to\mathrm{B}(G)_{\mathrm{basic}}$ 给出，其中 $e(J_\delta)\in\{\pm1\}$ 是 Kottitz 符号（[17]）\yizhuw{原文拼作“Kottitz”，应为 Kottwitz；译文照原拼写保留，以免与 [17] 的引用混淆。}。

这里必须对上述公式小心。
事实上，当 $J_\delta$ 模其中心是各向异性的时候，$x_\delta^*\cF_\varphi$ 是有限维的，迹才有意义。
而当不是这种情况时，必须把它解释为 $e(J_\delta)$ 乘以有限长表示 $x_\delta^*\cF_\varphi$ 在 $\delta$ 处的特征标。
更确切地说，这定义为 $\tfrac{1}{\mathrm{vol}(K\delta K)}\,\mathbf{1}_{K\delta K}$ 的迹，其中 $K$ 充分小。
当然，在第二种情形下这或许显得人为；但无论如何，这个稳定特征标能够被解释为 Frobenius 函数的迹，这一点让作者着迷\yizhu{当 $J_\delta$ 模中心紧时，$x_\delta^*\cF_\varphi$ 是有限维表示，可直接取迹；否则须用紧支函数 $\mathbf 1_{K\delta K}$ 的归一化迹来定义同一特征标。}{原文 \S5, p.30}。

\begin{remark}
作者曾试图把 Kottwitz 符号 $e(J_\delta)$ 纳入上同调平移或其他地方，却未能把它定成一个说得过去的东西。
前面已经出现过一次与 Frobenius 相联系的符号选择：当时我们说，当 $\mu$ 是 minuscule 时\yizhuw{minuscule 指余特征标 $\mu$ 在每条单根下的系数至多为 $1$；中文常译“微小”，此处保留英文。}，
\[
IC_\mu=\ol{\QQ}_\ell(\langle\rho,\mu\rangle)[\langle2\rho,\mu\rangle],
\]
而这个公式涉及选取 $\ol{\QQ}_\ell(1)$ 的一个平方根 $\ol{\QQ}_\ell(\tfrac12)$\yizhu{这里的 $\rho$ 是正根之和的一半（与前面表示标号的 $\rho$ 无关）：配对 $\langle\rho,\mu\rangle$ 给出扭转次数，$\langle2\rho,\mu\rangle$ 给出平移次数，故出现半整数次扭转。}{原文 \S5, 注 5.1, p.31}。
这一切似乎是一个微妙的问题。
\end{remark}

\begin{remark}
当然，若取 Frobenius 的迹乘以 $S_\varphi$ 的任意元素（而不是 $1\in S_\varphi$），就得到 $G$ 上的内视分布。
\end{remark}

\begin{remark}
设 $n\ge1$，$E_n|E$ 是 $E$ 的 $n$ 次非分歧扩域。
有映射 $G(E_n)/\sigma\text{-共轭}\to\mathrm{B}(G)$。
考虑元素 $\delta\in G(E_n)$，其范 $\delta\delta^\sigma\cdots\delta^{\sigma^{n-1}}$ 在 $G(E)$ 中的稳定共轭类为椭圆正则。
于是 $[\delta]\in\mathrm{B}(G)$ 是基本的，并有 $x_\delta:\Spa(\FF_{q^n})\to\Bun_G$。
可以把前述特征层性质加强为要求
\[
\mathrm{Tr}\bigl(\mathrm{Frob}_{q^n};x_\delta^*\cF_\varphi\bigr)
=(\mathrm{BC}_{E_n/E}T_\varphi)(\delta),
\]
即与 $\varphi$ 相联系的稳定特征标经基变换后在 $\delta$ 处取的值。
\end{remark}

\begin{remark}
特征层性质的灵感来自 [25] 中使用的方法：在 $p$-可除群的形变空间上取 Frobenius 的迹，以此在 $p$-adic 群上定义函数。
\end{remark}
